\section[Cote 62 : l'invariant $j$ d'une conique munie d'une trisection étale]{Cote 62 : l'invariant $j$ d'une conique munie d'une trisection étale\protect\verifie{Folder62}}\label{sec:62}

La cote 62, \emph{Courbes elliptiques (généralités -- vieille rédaction)}, notes que l'inventaire date \q{[à partir de 1963]}, décrit aux pages~41 à~53 les triplets $(X, X', s)$ formés d'une conique, d'une trisection étale et d'une section disjointe de $X'$, sur une base où $6$ est inversible. Dans une coordonnée qui envoie $X'$ sur $\{0, 1, \infty\}$, les six substitutions anharmoniques $t$, $1/t$, $1 - t$, $t/(t-1)$, $1/(1-t)$, $(t-1)/t$ forment $\mathfrak S_3$, et la page~39 calcule le quotient (11),
\[
  j(t) = -\frac{1}{27}\Bigl(\frac{(t+1)(2t-1)(t-2)}{t(t-1)}\Bigr)^2 ,
\]
revêtement galoisien de groupe $\mathfrak S_3$ non ramifié au-dessus de $\mathbf P^1 \smallsetminus \{0, 1, \infty\}$ ; l'involution $\kappa(t) = (t-2)/(2t-1)$ donne $j(\kappa t) = 1/j(t)$ (12) (\lot{62}{2}{39}). L'invariant du triplet est $j(s)$, et il détermine le triplet là où $j(j-1)$ est inversible (\lot{62}{3}{41}). Lean démontre cet énoncé fibre à fibre, sur un corps $K$ où $6$ est inversible, pour $t \notin \{0, 1\}$.

\begin{theoreme}[D'après les pages 39 à 47, l'invariant $j$\piste{62-triples-rigidified-moduli}]\label{thm:62-j}
Soit $K$ un corps où $6$ est inversible, et soient $s, t \in K \smallsetminus \{0, 1\}$. Alors $j(s) = j(t)$ si et seulement si $s$ est l'une des six images anharmoniques de $t$ ; en particulier $j$ est invariant par $\mathfrak S_3$. On a $j(t) = 0$ si et seulement si $t \in \{-1, 2, \frac12\}$, et $j(t) = 1$ si et seulement si $t^2 - t + 1 = 0$ ; aucune substitution non triviale ne fixe $t$ si et seulement si $j(t)(j(t) - 1) \neq 0$. Enfin $j = 1 - j_{\mathrm{cl}}/1728$, avec $j_{\mathrm{cl}}(t) = 256\,(t^2 - t + 1)^3/t^2(t-1)^2$, et $j(\kappa t)\,j(t) = 1$ pour $t \notin \{-1, 2, \frac12\}$.
\end{theoreme}

\begin{sloppypar}
\noindent L'énoncé tient tel que la lecture le donne (\texttt{j\_eq\_iff}, \texttt{stabilizer\_trivial\_iff}, \texttt{j\_eq\_one\_sub\_jcl}, \texttt{j\_kappa\_mul}, dans \href{\leanblob/Folder62.lean}{\texttt{Folder62.lean}}). La différence $j(s) - j(t)$ vaut $\frac{4}{27}$ fois le produit des six facteurs $s - g(t)$, dénominateurs chassés, divisé par $(s(s-1)\,t(t-1))^2$ ; l'invariance et les fibres s'en déduisent d'un coup. Le rapprochement avec $j_{\mathrm{cl}}$ est de l'édition.
\end{sloppypar}

\noindent La preuve vérifiée ne connaît ni coniques, ni trisections, ni base $S$. Le groupe $\mathfrak S_3$ et ses deux orbites exceptionnelles, $\{-1, 2, \frac12\}$ au-dessus de $j = 0$ et les racines de $t^2 - t + 1$ au-dessus de $j = 1$, sont ceux du cocycle anharmonique de la section~\ref{sec:146}, où les points fixes $-1$ et $\zeta$ donnent l'exposant~$6$ ; la droite $U = \mathbf P^1 \smallsetminus \{0, 1, \infty\}$ est le $U_{0,3}$ de l'épinglage de Legendre de la section~\ref{sec:64}. En termes de champs, $j$ est l'application de $U$ vers son quotient grossier par $\mathfrak S_3$, et les triplets $(X, X', s)$ sont les points du champ quotient $[U/\mathfrak S_3]$ ; cette formulation est de l'édition et n'est pas démontrée ici.

\begin{encadre}
\begin{theoreme*}[D'après les pages 41 à 53, trois descriptions des triplets]\nonverifie
Sur une base où $6$ est inversible, les triplets $(X, X', s)$ se décrivent de trois façons équivalentes : par un $\mathfrak S_3$-torseur $T$ muni d'une fonction $t$ sur $T$, $t$ et $t - 1$ inversibles, telle que $g \cdot t = g(t)$ ; par un faisceau inversible muni d'une trisection étale de somme nulle de son fibré vectoriel ; ou, par le quotient $q : X \to \mathbf P^1$ de degré~$6$, par l'invariant $j \in \Gamma(S, \mathcal O_S)$, qui détermine le triplet là où $j(j-1)$ est inversible.
\end{theoreme*}
\end{encadre}
{\small La transcription est très lacunaire ici : l'ouverture de la page~49 et la plus grande partie des pages~51 et~53 sont illisibles ou incertaines. Mathlib n'a pas les courbes de genre $0$ sur une base ni leurs trisections étales. On se reportera aux notes de Grothendieck (\lot{62}{3}{41}).\par}
